\section{Problema 4: Búsqueda en entornos parcialmente observables}

\subsection{Modelado del problema}
Explicar la situación en la que el robot no conoce exactamente su estado real.

\subsubsection{Estado}
Conjunto de estados posibles en vez de un único estado exacto.

\subsubsection{Percepción}
Lecturas de sensores, señales del entorno o información parcial.

\subsubsection{Objetivo}
Actuar correctamente aun con incertidumbre sobre el estado real.

\subsubsection{Dificultad principal}
El agente debe razonar sobre creencias y no sobre estados completamente conocidos.

\subsection{Algoritmo o enfoque elegido}
Estado de creencia, actualización por percepción y, si aplica, búsqueda en línea.

\subsection{Funcionamiento paso a paso}
Explicar cómo una acción y una observación reducen la incertidumbre.

\subsection{Limitaciones}
Ambigüedad, ruido en sensores y mayor dificultad de planificación.